\section*{Capítulo \ref{ch:g1:landscapes-around-us} --- Los paisajes que nos rodean}

\begin{solution}{exo:g1:landscapes-around-us:1}
\omterm{def:g1:landscapes-around-us:landform}{Montaña}, \omterm{def:g1:landscapes-around-us:landform}{colina}, \omterm{def:g1:landscapes-around-us:landform}{valle}, \omterm{def:g1:landscapes-around-us:landform}{llanura}.
\end{solution}

\begin{solution}{exo:g1:landscapes-around-us:2}
Una \omterm{def:g1:landscapes-around-us:landform}{montaña}: es más alta y más empinada que una \omterm{def:g1:landscapes-around-us:landform}{colina}.
\end{solution}

\begin{solution}{exo:g1:landscapes-around-us:3}
Naturales: el \omterm{def:g1:landscapes-around-us:water}{río}, el bosque, la playa. Construidos: la carretera, el
puente, la casa.
\end{solution}

\begin{solution}{exo:g1:landscapes-around-us:4}
Donde la tierra se encuentra con el mar.
\end{solution}

\begin{solution}{exo:g1:landscapes-around-us:5}
Por ejemplo: natural, las \omterm{def:g1:landscapes-around-us:landform}{montañas}, el bosque de las laderas o el arroyo;
construido, una casa o la iglesia del pueblo.
\end{solution}

\begin{solution}{exo:g1:landscapes-around-us:6}
La \omterm{def:g1:landscapes-around-us:landform}{montaña}, con nieve en la cima. El \omterm{def:g1:landscapes-around-us:water}{río} corre por el fondo del \omterm{def:g1:landscapes-around-us:landform}{valle},
entre las \omterm{def:g1:landscapes-around-us:landform}{montañas} y la \omterm{def:g1:landscapes-around-us:landform}{colina}.
\end{solution}

\begin{solution}{exo:g1:landscapes-around-us:7}
Un \omterm{def:g1:landscapes-around-us:water}{río} corre desde las tierras más altas (\omterm{def:g1:landscapes-around-us:landform}{colinas} y \omterm{def:g1:landscapes-around-us:landform}{montañas}) hasta un \omterm{def:g1:landscapes-around-us:water}{lago}
o hasta el mar. El agua siempre baja, así que no puede subir hasta lo alto
de una \omterm{def:g1:landscapes-around-us:landform}{colina}.
\end{solution}

\begin{solution}{exo:g1:landscapes-around-us:8}
El árbol está cerca de nosotros, en el primer plano; la \omterm{def:g1:landscapes-around-us:landform}{colina} está lejos,
en el fondo. Las cosas lejanas parecen más pequeñas que las cercanas.
\end{solution}

\begin{solution}{exo:g1:landscapes-around-us:9}
Primer plano: hierba alta y unos pocos arbustos en la parte de abajo de la
foto. Plano medio: el \omterm{def:g1:landscapes-around-us:water}{río} que serpentea por la \omterm{def:g1:landscapes-around-us:landform}{llanura} con árboles en sus
orillas, campos, una granja y una carretera.
Fondo: la tierra llana que sigue hasta el horizonte, y el cielo.
\end{solution}

\begin{solution}{exo:g1:landscapes-around-us:10}
Porque lo ha hecho la gente: desbrozó la tierra, la aró y sembró el trigo.
Sin la gente, allí habría hierba, arbustos o bosque, no un campo de trigo.
\end{solution}

\begin{solution}{exo:g1:landscapes-around-us:11}
Sin cambios: la \omterm{def:g1:landscapes-around-us:landform}{colina} y el \omterm{def:g1:landscapes-around-us:water}{río}. Añadido por la gente: los campos, más
casas, la carretera y el puente (y se ha talado casi todo el bosque).
\end{solution}
